\chapter{Metric Formulas}
\label{app:formulas}

This appendix gives the exact, implementation-level formula behind every
metric reported in \Cref{ch:results}, and behind a handful of further
diagnostic columns computed alongside them but not tabulated there (notably
Kyle's lambda). Each entry carries a scope tag: \textbf{U}niversal (both
engines), \textbf{C}DA-only, \textbf{F}BA-only, or \textbf{X} (defined but
always empty --- the dataset lacks what it would need). A metric marked
\textbf{X} is kept here rather than silently dropped, so that its absence
from every result in \Cref{ch:results} is a documented property of the
data, not an omission.

\section{Bucketing and Notation}

Every column is aggregated into fixed-width time buckets:
\begin{equation}
  \mathrm{bucket}(t) = \mathrm{anchor} + \left\lfloor \frac{t -
  \mathrm{anchor}}{\tau} \right\rfloor \tau,
\end{equation}
where $\mathrm{anchor}$ is the timestamp of the first record of the run
and $\tau$ is the interval (1\,second throughout this thesis,
\Cref{app:config}). Buckets are contiguous and empty-row-inclusive: a
bucket with no qualifying observation still produces a row, so a quiet
period is a gap in the data, not a gap in the file.

A missing value (\emph{None}, an empty CSV field) and an actual zero are
kept distinct and mean different things: \emph{None} is ``no qualifying
observation this bucket,'' zero is ``there were observations and the
value is genuinely zero.'' Every metric below follows this convention
with one deliberate exception, the intra-interval price dispersion
entry of \Cref{sec:formularq22} below. Price- and notional-denominated raw
values are stored in a fixed-point integer scale,
$\mathrm{PRICE\_SCALE} = 10^{6}$; bps and ratio metrics are already
unit-free and are \emph{not} rescaled.

\section{RQ2.1: Liquidity and Transaction Cost Metrics}

\paragraph{Quoted spread (U).} CDA: for every book snapshot in the
bucket with both sides present and $a_t + b_t > 0$,
$\mathrm{spread}_{\mathrm{bps}} = (a_t - b_t)/\midp \cdot 10^{4}$,
averaged over qualifying snapshots. FBA: for every batch with both a
best unfilled buy and a best unfilled sell and $\clr > 0$,
$(\text{unfilled sell} - \text{unfilled buy})/\clr \cdot 10^{4}$,
averaged over qualifying batches. A snapshot or batch that fails its own
gate is excluded from the mean, not counted as zero.

\paragraph{Depth at best (U).} CDA: mean of (best bid size $+$ best ask
size)$/2$ across snapshots --- touch only, not the whole book. FBA: mean
of (demand at $\clr$ $+$ supply at $\clr$)$/2$ across batches.

\paragraph{Depth within $x$ bps (U), $x \in \{10, 50, 100\}$.} For every
resting order, $\mathrm{bps\_away} = |p - \mathrm{ref}| \cdot
10^{4}/\mathrm{ref}$ against the snapshot midpoint or batch clearing
price; its quantity accumulates into every threshold with
$\mathrm{bps\_away} \le x$ --- the three columns are cumulative, so a
level within 10\,bps also counts toward the 50 and 100\,bps columns.

\paragraph{Book imbalance (C).} $(V^{bid} - V^{ask})/(V^{bid}+V^{ask})$
at the top of book, averaged over snapshots with nonzero touch depth.
Never set on an FBA row --- no code path computes an FBA analogue.

\paragraph{Total book depth (C).} Mean of (bid depth $+$ ask
depth)$/2$ across the \emph{whole} book, every price level both sides
--- contrast with depth at best's touch-only figure.

\paragraph{Effective spread (U), raw per-bucket value.}
\begin{equation}
  \mathrm{dev}_{\mathrm{bps}}(p_k, m_k, D_k) =
  \begin{cases}
    2(p_k - m_k)/m_k \cdot 10^{4} & D_k = \text{buy}\\
    -2(p_k - m_k)/m_k \cdot 10^{4} & D_k = \text{sell}\\
    2|p_k - m_k|/m_k \cdot 10^{4} & D_k \text{ unknown (FBA)}
  \end{cases}
\end{equation}
quantity-weighted per bucket, $m_k$ the pre-trade reference price (CDA:
book midpoint; FBA: the batch's own clearing price). \emph{None} if no
trade in the bucket had a usable reference price. The FBA has no
taker/maker distinction in a uniform-price batch, hence the third case:
every FBA contribution to this raw column is $\ge 0$ by construction (a
mean absolute deviation), while a CDA contribution can be negative (a
fill better than the reference price). This is the value each engine's
simulator writes to its own timeseries CSV; \Cref{ch:results} does not
report the FBA side of it directly, for the reason below.

\paragraph{Signed FBA spread: a CDA-tick proxy.} The raw FBA column above
has no aggressor sign, so before \Cref{ch:results} reports it, each bucket
is assigned a direction independently of the FBA's own reference price:
over the buckets where both engines traded, a Lee--Ready zero-tick rule
\citep{leeready1991} on \emph{consecutive CDA prices} classifies that
second's FBA batch as buy-side if the CDA ticked up, sell-side if it
ticked down, inheriting the last nonzero tick on a flat CDA price. The
signed value reported is this tick sign times the raw unsigned column
above. Ticking off the FBA's own clearing-price history instead would be
circular (that history \emph{is} $m_k$ above, so the sign it would give
exactly undoes $\mathrm{dev}_{\mathrm{bps}}$'s own sign flip); using the
CDA's independent price series avoids this. \Cref{sec:res21} discusses
what the tick rule can and cannot recover.

\paragraph{Realized spread (U), $\Delta \in \{1,5,30\}$\,s.} Same trades
and gate as effective spread, but against the reference price at
$t_k+\Delta$ instead of before the trade, found by a binary search over
the combined, time-sorted series of CDA midpoints and FBA clearing
prices (a book snapshot wins a tie against a batch at the identical
timestamp). \emph{None} for a horizon with no trade whose markout falls
within the retained window. \Cref{ch:results} reports the FBA side signed
by the same CDA-tick construction as the effective spread, one sign per
bucket shared across all three horizons.

\paragraph{Price impact (U), $\Delta \in \{1,5,30\}$\,s.} Effective
spread minus realized spread at the matching horizon --- the
adverse-selection component of the effective spread. \Cref{ch:results}
reports this as signed effective spread minus signed realized spread,
recomputed from the two signed terms rather than signing the raw,
already-differenced column, and assigned only when both signed terms are
computable for that bucket and horizon.

\paragraph{Amihud illiquidity (U).} From the same combined reference-price
series, $\mathrm{close}$ is the last price observed in a bucket, carried
forward across empty buckets. For a bucket with a close and a computable
prior close $>0$:
\begin{equation}
  \mathrm{ret} = \frac{\mathrm{close}-\mathrm{prev\_close}}{\mathrm{prev\_close}},
  \qquad
  \mathrm{amihud} = \frac{|\mathrm{ret}|}{\mathrm{dollar\_volume}} \cdot 10^{6}
\end{equation}
if the bucket traded nonzero volume, else \emph{None} --- a real price
move on zero volume reports \emph{None}, not $\infty$.

\section{RQ2.2: Price Discovery Metrics}
\label{sec:formularq22}

\paragraph{Realized volatility (U).} From the reference-price series,
per bucket, consecutive-pair returns $r_i = (p_{i+1}-p_i)/p_i$
(a pair with a non-positive base price is skipped); the reported value
is $\sqrt{\sum_i r_i^2}$ --- the root sum of squares, the standard
realized-variance estimator, deliberately \emph{not} a standard
deviation about the mean. A single return in the bucket reports
$|r|$ directly. \emph{None} with fewer than two usable points.

\paragraph{Matched-timestamp volatility (U).} Unlike the preceding
per-bucket estimator, this is not a field either simulator exports
directly: it is computed by the Python analysis pipeline
from both engines' \texttt{vwap}
columns jointly, after the fact. Build the intersection $\mathcal{T}$ of
buckets where both engines report a trade that second (\texttt{vwap}
non-\emph{None} on both sides); form consecutive-pair log-returns
$r^{e}_i = \ln\!\big(vwap^{e}_{t_{i+1}}/vwap^{e}_{t_i}\big)$ for
$e \in \{\text{FBA}, \text{CDA}\}$ over the same $t_i \in \mathcal{T}$ for
both engines (a pair with a non-positive price on either endpoint of
either engine is skipped); group by the calendar day of $t_{i+1}$, and
report $\sqrt{\mathrm{mean}_i(r^{e}_i)^2} \cdot 10^{4}$ per engine per
day, in bps: a root \emph{mean} of squares, not the root \emph{sum} the
per-bucket formula above uses, so a busier day is not conflated with a
more volatile one. A day with fewer than 5 matched returns is
\emph{None} for both engines. Because $\mathcal{T}$ requires a trade on
\emph{both} sides, this is conditional on joint trading activity, useful
for comparing the engines head to head but not a whole-month volatility
estimate for either alone.

\paragraph{Matched-timestamp volatility, intraday variant (U).} The same
construction, without the day-level grouping: each $r^{e}_i$ is reported
directly as $|r^{e}_i| \cdot 10^{4}$ bps, one observation per matched
return rather than one per calendar day, trading the daily version's
noise-averaging for full time resolution ($n$ in the hundreds of
thousands rather than around 31).

\paragraph{Intra-interval price dispersion (U).}
From bucketed trade prices, $\mathrm{disp} =
\mathrm{stddev}(\{p_k\})/\mathrm{mean}(\{p_k\}) \cdot 10^{4}$ (population
standard deviation). This is the one column that departs from the
\emph{None}-versus-zero convention stated in \Cref{app:formulas}'s first
section: it reports $0.0$, not \emph{None}, whenever the bucket had at
least one trade but the mean price was non-positive. \emph{None} only
when the bucket had no trades at all. Approximately zero for the FBA by
construction (one clearing price per batch): empirically, exactly $0.00$
in the large majority but not literally all FBA-side buckets
(\Cref{ch:results} reports the exact share for the run behind this
thesis) --- a bucket occasionally contains trades from more than one
batch clear, a rare edge case this appendix states rather than asserting
a stronger guarantee than the implementation actually gives.

\paragraph{Kyle's lambda (U, but built differently per engine).} Per
bucket, the OLS slope through the origin, $\lambda = \sum_i x_i
y_i/\sum_i x_i^2$, in bps per unit of the base asset (not
$\mathrm{PRICE\_SCALE}$-denominated). CDA: one observation per
\emph{sweep} (same-timestamp trades sharing a side and a pre-trade
midpoint, grouped in timestamp order); $x$ is the net signed executed
quantity of the sweep, $y = (m_{\mathrm{post}}-m_{\mathrm{pre}})/
m_{\mathrm{pre}} \cdot 10^{4}$ with $m_{\mathrm{post}}$ the midpoint
5\,s after the sweep --- a fixed horizon, independent of the three
realized-spread horizons above. FBA: one observation per priced batch,
walked in strictly ascending order; $x$ is the batch's net order flow
(submitted buy quantity minus submitted sell quantity, captured before
price selection), $y$ the percentage change in $\clr$ from the previous
priced batch. \emph{None} if $\sum_i x_i^2 = 0$ for the bucket. On this
dataset the CDA value is positive-skewed (buy-side sweeps lift the
midpoint, on average); the FBA value clusters near zero --- a
consequence of the price-selection rule described in
\Cref{sec:environment}, which anchors to the previous clearing price and
is close to impervious to a single batch's flow imbalance, not a defect
in the estimator. \Cref{ch:results} reports both.

\paragraph{Pricing error (X, always empty in this dataset).}
$|\clr - p^{\mathrm{ref}}|$ or $|\midp - p^{\mathrm{ref}}|$ against an
external reference price, in bps. No code path in this implementation ever
assigns this column a value: the December~2025 SOL order-status archive
used here carries no oracle or mark-price feed, and none is
reconstructed elsewhere in the pipeline. Every row of both timeseries
CSVs has this field empty, for the entire month, on both engines --- a
documented, structural gap rather than a result of zero.

\section{RQ2.3 and Engine-Control Metrics}

\paragraph{Executed volume, executed notional (U).} Plain running sums
over the trade stream: quantity, and quantity $\times$ price in raw
$\mathrm{PRICE\_SCALE}$ units (summed before any division).

\paragraph{VWAP (U).} $\mathrm{executed\_notional}/\mathrm{executed\_volume}$
if the bucket executed nonzero volume, else \emph{None}.

\paragraph{Trade count (U).} A plain count, not \emph{None}-valued,
incremented once per trade.

\paragraph{Fill rate (U).} Bucketed by \emph{each order's own submission
bucket}, not the trade's bucket: $\sum \mathrm{filled\_qty} / \sum
\mathrm{orig\_qty}$ over orders first seen in that bucket. \emph{None} if
no order was first seen in the bucket. Because this thesis's interval is
1\,second and most orders take longer than that to reach their eventual
fill state, this bucketing choice --- not a low overall fill propensity
--- is the main driver of a low average value; \Cref{ch:results} reports
it alongside average time to execution, below, for that reason.

\paragraph{Average time to execution (U).}
For every order with a first fill, $(\text{first\_fill\_ts} -
\text{first\_seen\_ts})/10^{9}$ seconds, bucketed by first-seen time;
mean per bucket, \emph{None} if no order first seen in the bucket ever
received a fill.

\paragraph{Trader surplus (U, a plain value, not \emph{None}-valued).}
Accumulated per trade and per side independently: the buy side (when its
limit price is known) contributes $\max(0, \pi - p_k) \cdot q_k$, the
sell side $\max(0, p_k - \pi) \cdot q_k$ --- clipped at zero, so an
adverse fill never contributes a negative amount. This is realized
price improvement against each side's own submitted limit, in raw
$\mathrm{PRICE\_SCALE}$ units.

\paragraph{Order size inflation (U).} Bucketed by first-seen time,
aggregated per (bucket, participant) as (submitted total, filled total);
a participant contributes $\mathrm{submitted}/\mathrm{filled}$ to the
bucket's mean only when both are strictly positive --- an order that
never filled at all is excluded, deliberately, rather than folded in
through a floor on the denominator, which would let never-traded orders
dominate the statistic rather than genuinely oversized ones.
\emph{None} if no participant in the bucket had a computable ratio.

\paragraph{Order-to-trade ratio (U).} $\mathrm{bucket\_message\_count} /
\mathrm{trade\_count}$; the numerator counts \emph{every} message,
accepted and rejected, so its population does not line up one-to-one
with fill rate's. \emph{None} if the bucket executed no trades.

\paragraph{Boundary concentration (F).} Share of order arrivals landing
in the final 10\% of an FBA batch's own window; a zero-width window
contributes nothing. Numerator and denominator accumulate across every
batch closing in the bucket before dividing. \emph{None} if the bucket
had zero total messages across its batches.

\paragraph{Throughput, clearing latency (U, wall-clock and
non-deterministic across runs and machines).} Throughput is the bucket's
message count divided by the summed measured compute time spent on
those messages. Clearing latency is that same summed compute time
divided by a count, where the count means two different things
depending on engine: \emph{per submission} for the CDA (one book
snapshot, one timing, per actionable record) but \emph{per batch clear}
for the FBA (one timing per batch) --- the same column name measures two
different granularities of operation depending on engine, a caveat that
matters when \Cref{ch:results} compares the two directly. Neither column
is a claim about market quality; both exist to rule out ``one engine is
just slower'' as an alternative explanation for a measured difference
elsewhere.

\paragraph{Unexecuted residual share (F).} Summed across every batch
closing in the bucket, using the $D(\cdot)$/$S(\cdot)$ notation of
\Cref{sec:fbanotation}: numerator $\sum |\,D(\clr) - S(\clr)\,|$,
denominator $\sum \max(D(\clr),S(\clr))$. This is the structural
one-sidedness of the schedule \emph{at the clearing price} within a
batch, not the complement of fill rate and not ``the share of submitted
volume left unfilled'' in the everyday sense --- a batch can match every
unit it possibly can ($\min(D(\clr),S(\clr))$ executed in full) and still
report a high value here if demand and supply were far apart in size at
$\clr$, which is common when $\tau$ is short and a batch contains few
orders. \emph{None} if the denominator is zero.

\section{Known Gaps}

The caveats behind pricing error, throughput/latency, price impact, and
Kyle's lambda are each stated where the metric itself is defined above,
not repeated here. Whether the simulated volume and price resemble what
the venue itself recorded is checked in \Cref{sec:testfidelity,sec:testprice},
not here; the one gap that check leaves open, the FBA's convergence to the
CDA as $\tau \to 0$, was not run, since it needs an interval sweep this
thesis does not include.
